
\paragraph{Momentum grid interpolation}

The usual method for determining interpolation coefficients $\delta
_{p,l+1/2,i+1,j+1/2}^{\left( 0\right) },$ $\delta _{p,l+1/2,i,j+1/2}^{\left(
0\right) },$ $\delta _{\xi ,l+1/2,i+1/2,j+1}^{\left( 0\right) }$ and $\delta
_{\xi ,l+1/2,i+1/2,j}^{\left( 0\right) }$ so that the distribution function $%
f_{0}^{\left( 0\right) }$ may be estimated on the full grid (flux grid) from
its value on the half grid (distribution grid) is to satisfy the constraint
that numerically, the stationary solution $\lim_{t\rightarrow +\infty %
}f_{0}^{\left( 0\right) }=f_{M}$ is an exact Maxwellian whithout any
external accelerating mechanism (RF wave or Ohmic electric field). In this
case, under the influence of collisional slowing-down and pitch-angle only,
coefficients $D_{p\xi }^{\left( 0\right) }=D_{\xi p}^{\left( 0\right) }=0,$
and 
\begin{eqnarray}
\left. S_{p}^{\left( 0\right) }\right| _{l+1/2,i+1,j+1/2}^{(\infty )}
&=&-D_{pp,l+1/2,i+1,j+1/2}^{\left( 0\right) }\frac{f_{0,l+1/2,i+3/2,j+1/2}^{%
\left( 0\right) (\infty )}-f_{0,l+1/2,i+1/2,j+1/2}^{\left( 0\right) (\infty %
)}}{\Delta p_{i+1}}  \nonumber \\
&&+F_{p,l+1/2,i+1,j+1/2}^{\left( 0\right) }f_{0,l+1/2,i+1,j+1/2}^{\left(
0\right) (\infty )}
\end{eqnarray}
\begin{eqnarray}
\left. S_{p}^{\left( 0\right) }\right| _{l+1/2,i,j+1/2}^{(\infty )}
&=&-D_{pp,l+1/2,i,j+1/2}^{\left( 0\right) }\frac{f_{0,l+1/2,i+1/2,j+1/2}^{%
\left( 0\right) (\infty )}-f_{0,l+1/2,i-1/2,j+1/2}^{\left( 0\right) (\infty %
)}}{\Delta p_{i}}  \nonumber \\
&&+F_{p,l+1/2,i,j+1/2}^{\left( 0\right) }f_{0,l+1/2,i,j+1/2}^{\left(
0\right) (\infty )}
\end{eqnarray}
\begin{eqnarray}
\left. S_{\xi }^{\left( 0\right) }\right| _{l+1/2,i+1/2,j+1}^{(\infty )}
&=&+D_{\xi \xi ,l+1/2,i+1/2,j+1}^{\left( 0\right) }\sqrt{1-\xi _{0,j+1}^{2}}%
\frac{f_{0,i+1/2,j+3/2}^{\left( 0\right) (\infty )}-f_{0,i+1/2,j+1/2}^{%
\left( 0\right) (\infty )}}{p_{i+1/2}\Delta \xi _{0,j+1}}  \nonumber \\
&&+F_{\xi ,l+1/2,i+1/2,j+1}^{\left( 0\right) }f_{0,l+1/2,i+1/2,j+1}^{\left(
0\right) (\infty )}
\end{eqnarray}
\begin{eqnarray}
\left. S_{\xi }^{\left( 0\right) }\right| _{l+1/2,i+1/2,j}^{(\infty )}
&=&+D_{\xi \xi ,l+1/2,i+1/2,j}^{\left( 0\right) }\sqrt{1-\xi _{0,j}^{2}}%
\frac{f_{0,i+1/2,j+1/2}^{\left( 0\right) (\infty )}-f_{0,i+1/2,j-1/2}^{%
\left( 0\right) (\infty )}}{p_{i+1/2}\Delta \xi _{0,j}}  \nonumber \\
&&+F_{\xi ,l+1/2,i+1/2,j}^{\left( 0\right) }f_{0,l+1/2,i+1/2,j}^{\left(
0\right) (\infty )}
\end{eqnarray}
or 
\begin{eqnarray}
\Delta p_{i+1}\left. S_{p}^{\left( 0\right) }\right| _{l+1/2,i+1,j+1/2}^{(%
\infty )} &=&-D_{pp,l+1/2,i+1,j+1/2}^{\left( 0\right) }\left(
f_{0,l+1/2,i+3/2,j+1/2}^{\left( 0\right) (\infty %
)}-f_{0,l+1/2,i+1/2,j+1/2}^{\left( 0\right) (\infty )}\right)   \nonumber \\
&&+\Delta p_{i+1}F_{p,l+1/2,i+1,j+1/2}^{\left( 0\right)
}f_{0,l+1/2,i+1,j+1/2}^{\left( 0\right) (\infty )}
\end{eqnarray}
\begin{eqnarray}
\Delta p_{i}\left. S_{p}^{\left( 0\right) }\right| _{l+1/2,i,j+1/2}^{(\infty %
)} &=&-D_{pp,l+1/2,i,j+1/2}^{\left( 0\right) }\left(
f_{0,l+1/2,i+1/2,j+1/2}^{\left( 0\right) (\infty %
)}-f_{0,l+1/2,i-1/2,j+1/2}^{\left( 0\right) (\infty )}\right)   \nonumber \\
&&+\Delta p_{i}F_{p,l+1/2,i,j+1/2}^{\left( 0\right)
}f_{0,l+1/2,i,j+1/2}^{\left( 0\right) (\infty )}
\end{eqnarray}
\begin{eqnarray}
p_{i+1/2}\Delta \xi _{0,j+1}\left. S_{\xi }^{\left( 0\right) }\right|
_{l+1/2,i+1/2,j+1}^{(\infty )} &=&+D_{\xi \xi ,l+1/2,i+1/2,j+1}^{\left(
0\right) }\sqrt{1-\xi _{0,j+1}^{2}}\times   \nonumber \\
&&\left( f_{0,l+1/2,i+1/2,j+3/2}^{\left( 0\right) (\infty %
)}-f_{0,l+1/2,i+1/2,j+1/2}^{\left( 0\right) (\infty )}\right)   \nonumber \\
&&+p_{i+1/2}\Delta \xi _{0,j+1}F_{\xi ,l+1/2,i+1/2,j+1}^{\left( 0\right)
}f_{0,l+1/2,i+1/2,j+1}^{\left( 0\right) (\infty )}
\end{eqnarray}
\begin{eqnarray}
p_{i+1/2}\Delta \xi _{0,j}\left. S_{\xi }^{\left( 0\right) }\right|
_{l+1/2,i+1/2,j}^{(\infty )} &=&+D_{\xi \xi ,l+1/2,i+1/2,j}^{\left( 0\right)
}\sqrt{1-\xi _{0,j}^{2}}\times   \nonumber \\
&&\left( f_{0,l+1/2,i+1/2,j+1/2}^{\left( 0\right) (\infty %
)}-f_{0,l+1/2,i+1/2,j-1/2}^{\left( 0\right) (\infty )}\right)   \nonumber \\
&&+p_{i+1/2}\Delta \xi _{0,j}F_{\xi ,l+1/2,i+1/2,j}^{\left( 0\right)
}f_{0,l+1/2,i+1/2,j}^{\left( 0\right) (\infty )}
\end{eqnarray}

Since this constraint corresponds to the relation 
\begin{equation}
\lim_{t\rightarrow +\infty }\frac{\partial f_{0}^{\left( 0\right) }}{%
\partial t}=-\nabla _{\mathbf{p}}\mathbf{S}_{\mathbf{p}}^{\left( 0\right) }=0
\end{equation}
one can deduce, 
\begin{eqnarray}
&&\Delta p_{i+1}F_{p,l+1/2,i+1,j+1/2}^{\left( 0\right)
}f_{0,l+1/2,i+1,j+1/2}^{\left( 0\right) (\infty )}  \nonumber \\
&=&D_{pp,l+1/2,i+1,j+1/2}^{\left( 0\right) }\left(
f_{0,l+1/2,i+3/2,j+1/2}^{\left( 0\right) (\infty %
)}-f_{0,l+1/2,i+1/2,j+1/2}^{\left( 0\right) (\infty )}\right) 
\end{eqnarray}
\begin{eqnarray}
&&\Delta p_{i}F_{p,l+1/2,i,j+1/2}^{\left( 0\right)
}f_{0,l+1/2,i,j+1/2}^{\left( 0\right) (\infty )}  \nonumber \\
&=&D_{pp,l+1/2,i,j+1/2}^{\left( 0\right) }\left(
f_{0,l+1/2,i+1/2,j+1/2}^{\left( 0\right) (\infty %
)}-f_{0,l+1/2,i-1/2,j+1/2}^{\left( 0\right) (\infty )}\right) 
\end{eqnarray}
\begin{eqnarray}
&&p_{i+1/2}\Delta \xi _{0,j+1}F_{\xi ,l+1/2,i+1/2,j+1}^{\left( 0\right)
}f_{0,l+1/2,i+1/2,j+1}^{\left( 0\right) (\infty )}  \nonumber \\
&=&-D_{\xi \xi ,l+1/2,i+1/2,j+1}^{\left( 0\right) }\sqrt{1-\xi _{0,j+1}^{2}}%
\left( f_{0,l+1/2,i+1/2,j+3/2}^{\left( 0\right) (\infty %
)}-f_{0,l+1/2,i+1/2,j+1/2}^{\left( 0\right) (\infty )}\right) 
\end{eqnarray}
\begin{eqnarray}
&&p_{i+1/2}\Delta \xi _{0,j}F_{\xi ,l+1/2,i+1/2,j}^{\left( 0\right)
}f_{0,l+1/2,i+1/2,j}^{\left( 0\right) (\infty )}  \nonumber \\
&=&-D_{\xi \xi ,l+1/2,i+1/2,j}^{\left( 0\right) }\sqrt{1-\xi _{0,j}^{2}}%
\left( f_{0,l+1/2,i+1/2,j+1/2}^{\left( 0\right) (\infty %
)}-f_{0,l+1/2,i+1/2,j-1/2}^{\left( 0\right) (\infty )}\right) 
\end{eqnarray}
and replacing the distribution functions on the full grid by their
interpolated values, according to the relations 
\begin{eqnarray}
f_{0,l+1/2,i+1,j+1/2}^{\left( 0\right) (k+1)} &=&\left( 1-\delta
_{p,l+1/2,i+1,j+1/2}^{\left( 0\right) }\right)
f_{0,l+1/2,i+3/2,j+1/2}^{\left( 0\right) (k+1)}  \nonumber \\
&&+\delta _{p,l+1/2,i+1,j+1/2}^{\left( 0\right)
}f_{0,l+1/2,i+1/2,j+1/2}^{\left( 0\right) (k+1)}
\end{eqnarray}
\begin{eqnarray}
f_{0,l+1/2,i,j+1/2}^{\left( 0\right) (k+1)} &=&\left( 1-\delta
_{p,l+1/2,i,j+1/2}^{\left( 0\right) }\right) f_{0,l+1/2,i+1/2,j+1/2}^{\left(
0\right) (k+1)}  \nonumber \\
&&+\delta _{p,l+1/2,i,j+1/2}^{\left( 0\right)
}f_{0,l+1/2,i-1/2,j+1/2}^{\left( 0\right) (k+1)}
\end{eqnarray}
and 
\begin{eqnarray}
f_{0,i+1/2,j+1}^{\left( 0\right) (k+1)} &=&\left( 1-\delta _{\xi
,l+1/2,i+1/2,j+1}^{\left( 0\right) }\right) f_{0,l+1/2,i+1/2,j+3/2}^{\left(
0\right) (k+1)}  \nonumber \\
&&+\delta _{\xi ,l+1/2,i+1/2,j+1}^{\left( 0\right)
}f_{0,l+1/2,i+1/2,j+1/2}^{\left( 0\right) (k+1)}
\end{eqnarray}
\begin{eqnarray}
f_{0,i+1/2,j}^{\left( 0\right) (k+1)} &=&\left( 1-\delta _{\xi
,l+1/2,i+1/2,j}^{\left( 0\right) }\right) f_{0,l+1/2,i+1/2,j+1/2}^{\left(
0\right) (k+1)}  \nonumber \\
&&+\delta _{\xi ,l+1/2,i+1/2,j}^{\left( 0\right)
}f_{0,l+1/2,i+1/2,j-1/2}^{\left( 0\right) (k+1)}
\end{eqnarray}
one finds, 
\begin{eqnarray}
&&\Delta p_{i+1}F_{p,l+1/2,i+1,j+1/2}^{\left( 0\right) }\left[ \left(
1-\delta _{p,l+1/2,i+1,j+1/2}^{\left( 0\right) }\right)
f_{0,l+1/2,i+3/2,j+1/2}^{\left( 0\right) (k+1)}\right.   \nonumber \\
&&\left. +\delta _{p,l+1/2,i+1,j+1/2}^{\left( 0\right)
}f_{0,l+1/2,i+1/2,j+1/2}^{\left( 0\right) (k+1)}\right]   \nonumber \\
&=&D_{pp,l+1/2,i+1,j+1/2}^{\left( 0\right) }\left(
f_{0,l+1/2,i+3/2,j+1/2}^{\left( 0\right) (\infty %
)}-f_{0,l+1/2,i+1/2,j+1/2}^{\left( 0\right) (\infty )}\right) 
\end{eqnarray}
\begin{eqnarray}
&&\Delta p_{i}F_{p,l+1/2,i,j+1/2}^{\left( 0\right) }\left[ \left( 1-\delta
_{p,l+1/2,i,j+1/2}^{\left( 0\right) }\right) f_{0,l+1/2,i+1/2,j+1/2}^{\left(
0\right) (k+1)}\right.   \nonumber \\
&&\left. +\delta _{p,l+1/2,i,j+1/2}^{\left( 0\right)
}f_{0,l+1/2,i-1/2,j+1/2}^{\left( 0\right) (k+1)}\right]   \nonumber \\
&=&D_{pp,l+1/2,i,j+1/2}^{\left( 0\right) }\left(
f_{0,l+1/2,i+1/2,j+1/2}^{\left( 0\right) (\infty %
)}-f_{0,l+1/2,i-1/2,j+1/2}^{\left( 0\right) (\infty )}\right) 
\end{eqnarray}
\begin{eqnarray}
&&p_{i+1/2}\Delta \xi _{0,j+1}F_{\xi ,l+1/2,i+1/2,j+1}^{\left( 0\right)
}\left[ \left( 1-\delta _{\xi ,l+1/2,i+1/2,j+1}^{\left( 0\right) }\right)
f_{0,l+1/2,i+1/2,j+3/2}^{\left( 0\right) (k+1)}\right.   \nonumber \\
&&\left. +\delta _{\xi ,l+1/2,i+1/2,j+1}^{\left( 0\right)
}f_{0,l+1/2,i+1/2,j+1/2}^{\left( 0\right) (k+1)}\right]   \nonumber \\
&=&-D_{\xi \xi ,l+1/2,i+1/2,j+1}^{\left( 0\right) }\sqrt{1-\xi _{0,j+1}^{2}}%
\left( f_{0,l+1/2,i+1/2,j+3/2}^{\left( 0\right) (\infty %
)}-f_{0,l+1/2,i+1/2,j+1/2}^{\left( 0\right) (\infty )}\right) 
\end{eqnarray}
\begin{eqnarray}
&&p_{i+1/2}\Delta \xi _{0,j}F_{\xi ,l+1/2,i+1/2,j}^{\left( 0\right) }\left[
\left( 1-\delta _{\xi ,l+1/2,i+1/2,j}^{\left( 0\right) }\right)
f_{0,l+1/2,i+1/2,j+1/2}^{\left( 0\right) (k+1)}\right.   \nonumber \\
&&\left. +\delta _{\xi ,l+1/2,i+1/2,j}^{\left( 0\right)
}f_{0,l+1/2,i+1/2,j-1/2}^{\left( 0\right) (k+1)}\right]   \nonumber \\
&=&-D_{\xi \xi ,l+1/2,i+1/2,j}^{\left( 0\right) }\sqrt{1-\xi _{0,j}^{2}}%
\left( f_{0,l+1/2,i+1/2,j+1/2}^{\left( 0\right) (\infty %
)}-f_{0,l+1/2,i+1/2,j-1/2}^{\left( 0\right) (\infty )}\right) 
\end{eqnarray}

Hence, the ratios are 
\begin{eqnarray}
&&\frac{f_{0,l+1/2,i+3/2,j+1/2}^{\left( 0\right) (\infty )}}{%
f_{0,l+1/2,i+1/2,j+1/2}^{\left( 0\right) (\infty )}}  \nonumber \\
&=&\frac{D_{pp,l+1/2,i+1,j+1/2}^{\left( 0\right) }+\Delta
p_{i+1}F_{p,l+1/2,i+1,j+1/2}^{\left( 0\right) }\delta
_{p,l+1/2,i+1,j+1/2}^{\left( 0\right) }}{D_{pp,l+1/2,i+1,j+1/2}^{\left(
0\right) }-\Delta p_{i+1}F_{p,l+1/2,i+1,j+1/2}^{\left( 0\right) }\left(
1-\delta _{p,l+1/2,i+1,j+1/2}^{\left( 0\right) }\right) }
\end{eqnarray}
\begin{eqnarray}
&&\frac{f_{0,l+1/2,i+1/2,j+1/2}^{\left( 0\right) (\infty )}}{%
f_{0,l+1/2,i-1/2,j+1/2}^{\left( 0\right) (\infty )}}  \nonumber \\
&=&\frac{D_{pp,l+1/2,i,j+1/2}^{\left( 0\right) }+\Delta
p_{i}F_{p,l+1/2,i,j+1/2}^{\left( 0\right) }\delta _{p,l+1/2,i,j+1/2}^{\left(
0\right) }}{D_{pp,l+1/2,i,j+1/2}^{\left( 0\right) }-\Delta
p_{i}F_{p,l+1/2,i,j+1/2}^{\left( 0\right) }\left( 1-\delta
_{p,l+1/2,i,j+1/2}^{\left( 0\right) }\right) }
\end{eqnarray}

\begin{eqnarray}
&&\frac{f_{0,l+1/2,i+1/2,j+3/2}^{\left( 0\right) (\infty )}}{%
f_{0,l+1/2,i+1/2,j+1/2}^{\left( 0\right) (\infty )}}  \nonumber \\
&=&\frac{-D_{\xi \xi ,l+1/2,i+1/2,j+1}^{\left( 0\right) }\sqrt{1-\xi
_{0,j+1}^{2}}+p_{i+1/2}\Delta \xi _{0,j+1}F_{\xi ,l+1/2,i+1/2,j+1}^{\left(
0\right) }\delta _{\xi ,l+1/2,i+1/2,j+1}^{\left( 0\right) }}{-D_{\xi \xi
,l+1/2,i+1/2,j+1}^{\left( 0\right) }\sqrt{1-\xi _{0,j+1}^{2}}%
-p_{i+1/2}\Delta \xi _{0,j+1}F_{\xi ,l+1/2,i+1/2,j+1}^{\left( 0\right)
}\left( 1-\delta _{\xi ,l+1/2,i+1/2,j+1}^{\left( 0\right) }\right) }
\end{eqnarray}

\begin{eqnarray}
&&\frac{f_{0,l+1/2,i+1/2,j+1/2}^{\left( 0\right) (\infty )}}{%
f_{0,l+1/2,i+1/2,j-1/2}^{\left( 0\right) (\infty )}}  \nonumber \\
&=&\frac{-D_{\xi \xi ,l+1/2,i+1/2,j}^{\left( 0\right) }\sqrt{1-\xi _{0,j}^{2}%
}+p_{i+1/2}\Delta \xi _{j}F_{\xi ,l+1/2,i+1/2,j}^{\left( 0\right) }\delta
_{\xi ,l+1/2,i+1/2,j}^{\left( 0\right) }}{-D_{\xi \xi
,l+1/2,i+1/2,j}^{\left( 0\right) }\sqrt{1-\xi _{0,j}^{2}}-p_{i+1/2}\Delta
\xi _{0,j}F_{\xi ,l+1/2,i+1/2,j}^{\left( 0\right) }\left( 1-\delta _{\xi
,l+1/2,i+1/2,j}^{\left( 0\right) }\right) }
\end{eqnarray}
which may be recast in a compact form 
\begin{equation}
\frac{f_{0,l+1/2,i+3/2,j+1/2}^{\left( 0\right) (\infty )}}{%
f_{0,l+1/2,i+1/2,j+1/2}^{\left( 0\right) (\infty )}}=\frac{%
1-w_{p,l+1/2,i+1,j+1/2}^{\left( 0\right) (u)}\delta
_{p,l+1/2,i+1,j+1/2}^{\left( 0\right) }}{1+w_{p,l+1/2,i+1,j+1/2}^{\left(
0\right) (u)}\left( 1-\delta _{p,l+1/2,i+1,j+1/2}^{\left( 0\right) }\right) }
\end{equation}
\begin{equation}
\frac{f_{0,l+1/2,i+1/2,j+1/2}^{\left( 0\right) (\infty )}}{%
f_{0,l+1/2,i-1/2,j+1/2}^{\left( 0\right) (\infty )}}=\frac{%
1-w_{p,i,j+1/2}^{\left( 0\right) (u)}\delta _{p,l+1/2,i,j+1/2}^{\left(
0\right) }}{1+w_{p,l+1/2,i,j+1/2}^{\left( 0\right) (u)}\left( 1-\delta
_{p,l+1/2,i,j+1/2}^{\left( 0\right) }\right) }
\end{equation}

\begin{equation}
\frac{f_{0,l+1/2,i+1/2,j+3/2}^{\left( 0\right) (\infty )}}{%
f_{0,l+1/2,i+1/2,j+1/2}^{\left( 0\right) (\infty )}}=\frac{1-w_{\xi
,l+1/2,i+1/2,j+1}^{\left( 0\right) (u)}\delta _{\xi
,l+1/2,i+1/2,j+1}^{\left( 0\right) }}{1+w_{\xi ,l+1/2,i+1/2,j+1}^{(u)}\left(
1-\delta _{\xi ,l+1/2,i+1/2,j+1}^{\left( 0\right) }\right) }
\end{equation}

\begin{equation}
\frac{f_{0,l+1/2,i+1/2,j+1/2}^{\left( 0\right) (\infty )}}{%
f_{0,,l+1/2,i+1/2,j-1/2}^{\left( 0\right) (\infty )}}=\frac{1-w_{\xi
,l+1/2,i+1/2,j}^{\left( 0\right) (u)}\delta _{\xi ,l+1/2,i+1/2,j}^{\left(
0\right) }}{1+w_{\xi ,l+1/2,i+1/2,j}^{^{\left( 0\right) }(u)}\left( 1-\delta
_{\xi ,l+1/2,i+1/2,j}^{\left( 0\right) }\right) }
\end{equation}
with 
\begin{equation}
w_{p,l+1/2,i+1,j+1/2}^{\left( 0\right) (u)}=-\Delta p_{i+1}\frac{
F_{p,l+1/2,i+1,j+1/2}^{\left( 0\right) }}{D_{pp,l+1/2,i+1,j+1/2}^{\left(
0\right) }}
\end{equation}
\begin{equation}
w_{p,l+1/2,i,j+1/2}^{\left( 0\right) (u)}=-\Delta p_{i}\frac{
F_{p,l+1/2,i,j+1/2}^{\left( 0\right) }}{D_{pp,l+1/2,i,j+1/2}^{\left(
0\right) }}
\end{equation}
\begin{equation}
w_{\xi ,l+1/2,i+1/2,j+1}^{\left( 0\right) (u)}=+p_{i+1/2}\frac{\Delta \xi
_{0,j+1}}{\sqrt{1-\xi _{0,j+1}^{2}}}\frac{F_{\xi ,l+1/2,i+1/2,j+1}^{\left(
0\right) }}{D_{\xi \xi ,l+1/2,i+1/2,j+1}^{\left( 0\right) }}
\end{equation}
\begin{equation}
w_{\xi ,l+1/2,i+1/2,j}^{\left( 0\right) (u)}=+p_{i+1/2}\frac{\Delta \xi
_{0,j}}{\sqrt{1-\xi _{0,j}^{2}}}\frac{F_{\xi ,l+1/2,i+1/2,j}^{\left(
0\right) }}{D_{\xi \xi ,l+1/2,i/2,j}^{\left( 0\right) }}
\end{equation}

Furthermore, since $\nabla _{\mathbf{p}}\mathbf{S}_{\mathbf{p}}^{\left(
0\right) }=0,$ one deduces naturally that 
\begin{equation}
-D_{pp}^{\left( 0\right) }\frac{df_{0}^{\left( 0\right) (\infty )}}{dp}%
+F_{p}^{\left( 0\right) }f_{0}^{\left( 0\right) (\infty )}=0
\end{equation}
\begin{equation}
+\frac{\sqrt{1-\xi _{0}^{2}}}{p}D_{\xi \xi }^{\left( 0\right) }\frac{%
df_{0}^{\left( 0\right) (\infty )}}{d\xi _{0}}+F_{\xi }^{\left( 0\right)
}f_{0}^{\left( 0\right) (\infty )}=0
\end{equation}
and since $F_{\xi }^{\left( 0\right) }=0$ for collisions, because of the
isotropic nature of pitch-angle scattering 
\begin{equation}
\frac{\partial f_{0}^{\left( 0\right) (\infty )}}{f_{0}^{(\infty )}}=\left( 
\frac{F_{p}^{\left( 0\right) }}{D_{pp}^{\left( 0\right) }}\right) \partial p
\label{gridInterpLim}
\end{equation}
\begin{equation}
\frac{\partial f_{0}^{\left( 0\right) (\infty )}}{\partial \xi _{0}}=0
\end{equation}
which indicates that when collisions is the single physical process
considered, $f_{0}^{\left( 0\right) (\infty )}$ is independent of $\xi _{0}.$
Therefore, $\frac{\partial }{\partial p}\rightarrow \frac{d}{dp},$ and
integrating the first equation of (\ref{gridInterpLim}) leads to the
relation 
\begin{equation}
\ln f_{0}^{\left( 0\right) }=\int \left( \frac{F_{p}^{\left( 0\right) }}{%
D_{pp}^{\left( 0\right) }}\right) dp+Cte
\end{equation}

Values of $f_{0}^{\left( 0\right) (\infty )}$ on the half grid positions $%
\left( l+1/2,i+3/2,j+1/2\right) $ and $\left( l+1/2,i+1/2,j+1/2\right) $ are
then

\begin{equation}
\ln f_{0,l+1/2,i+3/2,j+1/2}^{\left( 0\right) (\infty
)}=\int_{0}^{p_{i+3/2}}\left( \frac{F_{p}^{\left( 0\right) }\left( \psi
_{l+1/2,}p,\xi _{0,j+1/2}\right) }{D_{pp}^{\left( 0\right) }\left( \psi
_{l+1/2,}p,\xi _{0,j+1/2}\right) }\right) dp+Cte
\end{equation}
\begin{equation}
\ln f_{0,l+1/2,i+1/2,j+1/2}^{\left( 0\right) (\infty
)}=\int_{0}^{p_{i+1/2}}\left( \frac{F_{p}^{\left( 0\right) }\left( \psi
_{l+1/2,}p,\xi _{0,j+1/2}\right) }{D_{pp}^{\left( 0\right) }\left( \psi
_{l+1/2,}p,\xi _{0,j+1/2}\right) }\right) dp+Cte
\end{equation}
and ratios are finally easily determined 
\begin{equation}
\frac{f_{0,l+1/2,i+3/2,j+1/2}^{\left( 0\right) (\infty )}}{%
f_{0,l+1/2,i+1/2,j+1/2}^{\left( 0\right) (\infty )}}=\exp \left[
\int_{p_{i+1/2}}^{p_{i+3/2}}\left( \frac{F_{p}^{\left( 0\right) }\left( \psi
_{l+1/2,}p,\xi _{0,j+1/2}\right) }{D_{pp}^{\left( 0\right) }\left( \psi
_{l+1/2,}p,\xi _{0,j+1/2}\right) }\right) dp\right]
\end{equation}
\begin{equation}
\frac{f_{0,l+1/2,i+1/2,j+1/2}^{\left( 0\right) (\infty )}}{%
f_{0,l+1/2,i-1/2,j+1/2}^{\left( 0\right) (\infty )}}=\exp \left[
\int_{p_{i-1/2}}^{p_{i13/2}}\left( \frac{F_{p}^{\left( 0\right) }\left( \psi
_{l+1/2,}p,\xi _{0,j+1/2}\right) }{D_{pp}^{\left( 0\right) }\left( \psi
_{l+1/2,}p,\xi _{0,j+1/2}\right) }\right) dp\right]
\end{equation}
while, since $f_{0}^{\left( 0\right) (\infty )}$ does not depend of $\xi
_{0},$ 
\begin{equation}
\frac{f_{0,l+1/2,i+1/2,j+3/2}^{\left( 0\right) (\infty )}}{%
f_{0,l+1/2,i+1/2,j+1/2}^{\left( 0\right) (\infty )}}=\frac{%
f_{0,l+1/2,i+1/2,j+1/2}^{\left( 0\right) (\infty )}}{%
f_{0,l+1/2,i+1/2,j-1/2}^{\left( 0\right) (\infty )}}=1
\end{equation}

It is then straightforward to evaluate coefficients $\delta
_{p,l+1/2,i+1,j+1/2}^{\left( 0\right) }$ et $\delta
_{p,l+1/2,i,j+1/2}^{\left( 0\right) }$ by the trapezoidal method, 
\begin{eqnarray}
&&\int_{p_{i+1/2}}^{p_{i+3/2}}\frac{F_{p}^{\left( 0\right) }\left( \psi
_{l+1/2,}p,\xi _{0,j+1/2}\right) }{D_{pp}^{\left( 0\right) }\left( \psi
_{l+1/2,}p,\xi _{0,j+1/2}\right) }dp\simeq   \nonumber \\
&&\frac{1}{2}\frac{F_{p,l+1/2,i+1/2,j+1/2}^{\left( 0\right) }}{%
D_{pp,l+1/2,i+1/2,j+1/2}^{\left( 0\right) }}\frac{\Delta p_{i+1/2}}{2} 
\nonumber \\
&&+\frac{1}{2}\frac{F_{p,l+1/2,i+1,j+1/2}^{\left( 0\right) }}{%
D_{pp,l+1/2,i+1,j+1/2}^{\left( 0\right) }}\frac{\Delta p_{i+1/2}}{2} 
\nonumber \\
&&+\frac{1}{2}\frac{F_{p,l+1/2,i+1,j+1/2}^{\left( 0\right) }}{%
D_{pp,l+1/2,i+1,j+1/2}^{\left( 0\right) }}\frac{\Delta p_{i+3/2}}{2} 
\nonumber \\
&&+\frac{1}{2}\frac{F_{p,l+1/2,i+3/2,j+1/2}^{\left( 0\right) }}{%
D_{pp,l+1/2,i+3/2,j+1/2}^{\left( 0\right) }}\frac{\Delta p_{i+3/2}}{2}
\end{eqnarray}
and 
\begin{eqnarray}
&&\int_{p_{i-1/2}}^{p_{i+1/2}}\frac{F_{p}^{\left( 0\right) }\left( \psi
_{l+1/2,}p,\xi _{0,j+1/2}\right) }{D_{pp}^{\left( 0\right) }\left( \psi
_{l+1/2,}p,\xi _{0,j+1/2}\right) }dp\simeq   \nonumber \\
&&\frac{1}{2}\frac{F_{p,l+1/2,i-1/2,j+1/2}^{\left( 0\right) }}{%
D_{pp,l+1/2,i-1/2,j+1/2}^{\left( 0\right) }}\frac{\Delta p_{i-1/2}}{2} 
\nonumber \\
&&+\frac{1}{2}\frac{F_{p,l+1/2,i+3/2,j+1/2}^{\left( 0\right) }}{%
D_{pp,l+1/2,i+3/2,j+1/2}^{\left( 0\right) }}\frac{\Delta p_{i+3/2}}{2} 
\nonumber \\
&&+\frac{1}{2}\frac{F_{p,l+1/2,i,j+1/2}^{\left( 0\right) }}{%
D_{pp,l+1/2,i,j+1/2}^{\left( 0\right) }}\frac{\Delta p_{i+1/2}}{2}  \nonumber
\\
&&+\frac{1}{2}\frac{F_{p,l+1/2,i+1/2,j+1/2}^{\left( 0\right) }}{%
D_{pp,l+1/2,i+1/2,j+1/2}^{\left( 0\right) }}\frac{\Delta p_{i+1/2}}{2}
\end{eqnarray}
then gathering terms $p_{i+3/2},$ $p_{i+1}$ and $p_{i+1/2}$, one obtains 
\begin{eqnarray}
&&\int_{p_{i+1/2}}^{p_{i+3/2}}\frac{F_{p}^{\left( 0\right) }\left( \psi
_{l+1/2,}p,\xi _{0,j+1/2}\right) }{D_{pp}^{\left( 0\right) }\left( \psi
_{l+1/2,}p,\xi _{0,j+1/2}\right) }dp\simeq   \nonumber \\
&&\frac{1}{2}\frac{F_{p,l+1/2,i+1,j+1/2}^{\left( 0\right) }}{%
D_{pp,l+1/2,i+1,j+1/2}^{\left( 0\right) }}\left( \frac{\Delta p_{i+1/2}}{2}+%
\frac{\Delta p_{i+3/2}}{2}\right)   \nonumber \\
&&+\frac{F_{p,l+1/2,i+1/2,j+1/2}^{\left( 0\right) }}{%
D_{pp,l+1/2,i+1/2,j+1/2}^{\left( 0\right) }}\frac{\Delta p_{i+1/2}}{4} 
\nonumber \\
&&+\frac{F_{p,l+1/2,i+3/2,j+1/2}^{\left( 0\right) }}{%
D_{pp,l+1/2,i+3/2,j+1/2}^{\left( 0\right) }}\frac{\Delta p_{i+3/2}}{4} 
\nonumber \\
&=&\frac{F_{p,l+1/2,i+1,j+1/2}^{\left( 0\right) }}{D_{pp,l+1/2,i+1,j+1/2}^{%
\left( 0\right) }}\frac{\Delta p_{i+1}}{2}  \nonumber \\
&&+\frac{F_{p,l+1/2,i+1/2,j+1/2}^{\left( 0\right) }}{%
D_{pp,l+1/2,i+1/2,j+1/2}^{\left( 0\right) }}\frac{\Delta p_{i+1/2}}{4} 
\nonumber \\
&&+\frac{F_{p,l+1/2,i+3/2,j+1/2}^{\left( 0\right) }}{%
D_{pp,l+1/2,i+3/2,j+1/2}^{\left( 0\right) }}\frac{\Delta p_{i+3/2}}{4}
\end{eqnarray}
or 
\begin{eqnarray}
&&\int_{p_{i+1/2}}^{p_{i+3/2}}\frac{F_{p}^{\left( 0\right) }\left( \psi
_{l+1/2,}p,\xi _{0,j+1/2}\right) }{D_{pp}^{\left( 0\right) }\left( \psi
_{l+1/2,}p,\xi _{0,j+1/2}\right) }dp=  \nonumber \\
&&\Delta p_{i+1}\frac{F_{p,l+1/2,i+1,j+1/2}^{\left( 0\right) }}{%
D_{pp,l+1/2,i+1,j+1/2}^{\left( 0\right) }}  \nonumber \\
&&+\frac{F_{p,l+1/2,i+1/2,j+1/2}^{\left( 0\right) }}{%
D_{pp,l+1/2,i+1/2,j+1/2}^{\left( 0\right) }}\frac{\Delta p_{i+1/2}}{4} 
\nonumber \\
&&+\frac{F_{p,l+1/2,i+3/2,j+1/2}^{\left( 0\right) }}{%
D_{pp,l+1/2,i+3/2,j+1/2}^{\left( 0\right) }}\frac{\Delta p_{i+3/2}}{4} 
\nonumber \\
&&-\frac{\Delta p_{i+1}}{2}\frac{F_{p,l+1/2,i+1,j+1/2}^{\left( 0\right) }}{%
D_{pp,l+1/2,i+1,j+1/2}^{\left( 0\right) }}
\end{eqnarray}

Since 
\begin{eqnarray}
&&\frac{F_{p,l+1/2,i+1,j+1/2}^{\left( 0\right) }}{D_{pp,l+1/2,i+1,j+1/2}^{%
\left( 0\right) }}=  \nonumber \\
&&\frac{\Delta p_{i+3/2}}{\Delta p_{i+1/2}+\Delta p_{i+3/2}}\frac{%
F_{p,l+1/2,i+1/2,j+1/2}^{\left( 0\right) }}{D_{pp,l+1/2,i+1/2,j+1/2}^{\left(
0\right) }}  \nonumber \\
&&+\frac{\Delta p_{i+1/2}}{\Delta p_{i+1/2}+\Delta p_{i+3/2}}\frac{%
F_{p,l+1/2,i+3/2,j+1/2}^{\left( 0\right) }}{D_{pp,l+1/2,i+3/2,j+1/2}^{\left(
0\right) }}
\end{eqnarray}
one obtains 
\begin{eqnarray}
&&\frac{F_{p,l+1/2,i+1,j+1/2}^{\left( 0\right) }}{D_{pp,l+1/2,i+1,j+1/2}^{%
\left( 0\right) }}=  \nonumber \\
&&\frac{\Delta p_{i+3/2}}{2\Delta p_{i+1}}\frac{F_{p,l+1/2,i+1/2,j+1/2}^{%
\left( 0\right) }}{D_{pp,l+1/2,i+1/2,j+1/2}^{\left( 0\right) }}  \nonumber \\
&&+\frac{\Delta p_{i+1/2}}{2\Delta p_{i+1}}\frac{F_{p,l+1/2,i+3/2,j+1/2}^{%
\left( 0\right) }}{D_{pp,l+1/2,i+3/2,j+1/2}^{\left( 0\right) }}
\end{eqnarray}
and after gathering all terms 
\begin{eqnarray}
&&\int_{p_{i+1/2}}^{p_{i+3/2}}\frac{F_{p}^{\left( 0\right) }\left( \psi
_{l+1/2,}p,\xi _{0,j+1/2}\right) }{D_{pp}^{\left( 0\right) }\left( \psi
_{l+1/2,}p,\xi _{0,j+1/2}\right) }dp=  \nonumber \\
&&\Delta p_{i+1}\frac{F_{p,l+1/2,i+1,j+1/2}^{\left( 0\right) }}{%
D_{pp,l+1/2,i+1,j+1/2}^{\left( 0\right) }}  \nonumber \\
&&+\frac{\Delta p_{i+3/2}-\Delta p_{i+1/2}}{4}\frac{%
F_{p,l+1/2,i+3/2,j+1/2}^{\left( 0\right) }}{D_{pp,l+1/2,i+3/2,j+1/2}^{\left(
0\right) }}  \nonumber \\
&&-\frac{\Delta p_{i+3/2}-\Delta p_{i+1/2}}{4}\frac{%
F_{p,l+1/2,i+1/2,j+1/2}^{\left( 0\right) }}{D_{pp,l+1/2,i+1/2,j+1/2}^{\left(
0\right) }}
\end{eqnarray}

Much as the same way, using index conversion $p_{i+3/2}\rightarrow
p_{i+1/2}, $ $p_{i+1}\rightarrow p_{i}$ and $p_{i+1/2}\rightarrow p_{i-1/2}$
in the previous relation, one obtains for the second integral,

\begin{eqnarray}
&&\int_{p_{i-1/2}}^{p_{i+1/2}}\frac{F_{p}^{\left( 0\right) }\left( \psi
_{l+1/2,}p,\xi _{0,j+1/2}\right) }{D_{pp}^{\left( 0\right) }\left( \psi
_{l+1/2,}p,\xi _{0,j+1/2}\right) }dp\simeq   \nonumber \\
&&\Delta p_{i}\frac{F_{p,l+1/2,i,j+1/2}^{\left( 0\right) }}{%
D_{pp,l+1/2,i,j+1/2}^{\left( 0\right) }}  \nonumber \\
&&+\frac{\Delta p_{i+1/2}-\Delta p_{i-1/2}}{4}\frac{%
F_{p,l+1/2,i+1/2,j+1/2}^{\left( 0\right) }}{D_{pp,l+1/2,i+1/2,j+1/2}^{\left(
0\right) }}  \nonumber \\
&&-\frac{\Delta p_{i+1/2}-\Delta p_{i-1/2}}{4}\frac{%
F_{p,l+1/2,i-1/2,j+1/2}^{\left( 0\right) }}{D_{pp,l+1/2,i-1/2,j+1/2}^{\left(
0\right) }}
\end{eqnarray}

From the two formulations of the ratios 
\begin{equation}
\frac{f_{0,l+1/2,i+3/2,j+1/2}^{\left( 0\right) (\infty )}}{%
f_{0,l+1/2,i+1/2,j+1/2}^{\left( 0\right) (\infty )}}
\end{equation}
one obtains easily 
\begin{eqnarray}
\delta _{p,l+1/2,i+1,j+1/2}^{\left( 0\right) } &=&\frac{1}{%
w_{p,l+1/2,i+1,j+1/2}^{\left( 0\right) (u)}}  \nonumber \\
&&-\frac{1}{\exp \left[ w_{p,l+1/2,i+1,j+1/2}^{\left( 0\right)
(u)}+w_{p,l+1/2,i+1,j+1/2}^{\left( 0\right) (nu)}\right] -1}
\end{eqnarray}
with 
\begin{eqnarray}
w_{p,l+1/2,i+1,j+1/2}^{\left( 0\right) (nu)} &=&-\frac{\Delta p_{i+3/2}-%
\Delta p_{i+1/2}}{4}\frac{F_{p,l+1/2,i+3/2,j+1/2}^{\left( 0\right) }}{%
D_{pp,l+1/2,i+3/2,j+1/2}^{\left( 0\right) }}  \nonumber \\
&&+\frac{\Delta p_{i+3/2}-\Delta p_{i+1/2}}{4}\frac{%
F_{p,l+1/2,i+1/2,j+1/2}^{\left( 0\right) }}{D_{pp,l+1/2,i+1/2,j+1/2}^{\left(
0\right) }}
\end{eqnarray}

Much in the same way, 
\begin{equation}
\delta _{p,l+1/2,i,j+1/2}^{\left( 0\right) }=\frac{1}{w_{p,l+1/2,i,j+1/2}^{%
\left( 0\right) (u)}}-\frac{1}{\exp \left[ w_{p,l+1/2,i,j+1/2}^{\left(
0\right) (u)}+w_{p,l+1/2,i,j+1/2}^{\left( 0\right) (nu)}\right] -1}
\end{equation}
with \ss 
\begin{eqnarray}
w_{p,l+1/2,i,j+1/2}^{\left( 0\right) (nu)} &=&-\frac{\Delta p_{i+1/2}-\Delta %
p_{i-1/2}}{4}\frac{F_{p,l+1/2,i+1/2,j+1/2}^{\left( 0\right) }}{%
D_{pp,l+1/2,i+1/2,j+1/2}^{\left( 0\right) }}  \nonumber \\
&&+\frac{\Delta p_{i+1/2}-\Delta p_{i-1/2}}{4}\frac{%
F_{p,l+1/2,i-1/2,j+1/2}^{\left( 0\right) }}{D_{pp,l+1/2,i-1/2,j+1/2}^{\left(
0\right) }}
\end{eqnarray}

For a uniform grid, $w_{p,l+1/2,i,j+1/2}^{\left( 0\right) (nu)}=0,$ and the
well known expression of the Chang and Cooper function is recovered, 
\begin{equation}
g\left( x\right) =\frac{1}{x}-\frac{1}{e^{x}-1}
\end{equation}
with $x=w_{p}^{\left( 0\right) (u)}.$ In the limit where $x\ll 1,$ an
expansion of the function gives $g\left( x\right) =0.5-x/12+x^{3}/720+%
\mathrm{O}\left( x^{4}\right) $ with $\lim_{x\rightarrow 0}g\left( x\right) =%
\frac{1}{2}.$ For a non-uniform grid, a generalized Chang and Cooper
function must be introduced, with two arguments 
\begin{equation}
g\left( x,y\right) =\frac{1}{x}-\frac{1}{e^{x+y}-1}
\end{equation}
where the term $y=w_{p}^{\left( 0\right) (nu)}$ is zero for the uniform case.

With this definition, 
\begin{equation}
\frac{f_{0,l+1/2,i+3/2,j+1/2}^{\left( 0\right) (\infty )}}{%
f_{0,l+1/2,i+1/2,j+1/2}^{(\infty )}}=\frac{T_{1}^{cc}}{T_{2}^{cc}}=\frac{%
T_{4}^{cc}}{T_{5}^{cc}}
\end{equation}
where 
\begin{eqnarray}
T_{1}^{cc} &=&1-w_{p,l+1/2,i+1,j+1/2}^{\left( 0\right) (u)}\left[ \frac{1}{%
w_{p,,l+1/2,i+1,j+1/2}^{\left( 0\right) (u)}}\right.   \nonumber \\
&&\left. -\frac{1}{\exp \left[ w_{p,l+1/2,i+1,j+1/2}^{\left( 0\right)
(u)}+w_{p,l+1/2,i+1,j+1/2}^{\left( 0\right) (nu)}\right] -1}\right] 
\end{eqnarray}
\begin{equation}
T_{2}^{cc}=1+w_{p,l+1/2,i+1,j+1/2}^{\left( 0\right) (u)}\left(
1-T_{3}^{cc}\right) 
\end{equation}
and 
\begin{equation}
T_{3}^{cc}=\frac{1}{w_{p,l+1/2,i+1,j+1/2}^{\left( 0\right) (u)}}-\frac{1}{%
\exp \left[ w_{p,l+1/2,i+1,j+1/2}^{\left( 0\right)
(u)}+w_{p,l+1/2,i+1,j+1/2}^{\left( 0\right) (nu)}\right] -1}
\end{equation}
whilewith 
\begin{equation}
T_{4}^{cc}=\frac{w_{p,l+1/2,i+1,j+1/2}^{\left( 0\right) (u)}}{\exp \left[
w_{p,l+1/2,i+1,j+1/2}^{\left( 0\right) (u)}+w_{p,l+1/2,i+1,j+1/2}^{\left(
0\right) (nu)}\right] -1}
\end{equation}
and 
\begin{equation}
T_{5}^{cc}=w_{p,l+1/2,i+1,j+1/2}^{\left( 0\right) (u)}+\frac{%
w_{p,l+1/2,i+1,j+1/2}^{\left( 0\right) (u)}}{\exp \left[
w_{p,l+1/2,i+1,j+1/2}^{\left( 0\right) (u)}+w_{p,l+1/2,i+1,j+1/2}^{\left(
0\right) (nu)}\right] -1}
\end{equation}

This expression simplifies and becomes 
\begin{equation}
\frac{f_{0,l+1/2,i+3/2,j+1/2}^{\left( 0\right) (\infty )}}{
f_{0,l+1/2,i+1/2,j+1/2}^{\left( 0\right) (\infty )}}=\exp \left[
-w_{p,l+1/2,i+1,j+1/2}^{\left( 0\right) (u)}-w_{p,l+1/2,i+1,j+1/2}^{\left(
0\right) (nu)}\right]
\end{equation}

Furthermore, when plasma relaxation by collisions is the single process
considered, the relation 
\begin{equation}
\frac{F_{p}^{\left( 0\right) }\left( \psi ,p,\xi _{0}\right) }{%
D_{pp}^{\left( 0\right) }\left( \psi ,p,\xi _{0}\right) }\cong -p
\end{equation}
holds, since $f_{0}^{\left( 0\right) (\infty )}$ must be Maxwellian $f_{M},$
and one finds

\begin{equation}
\frac{f_{0,l+1/2,i+3/2,j+1/2}^{\left( 0\right) (\infty )}}{%
f_{0,l+1/2,i+1/2,j+1/2}^{\left( 0\right) (\infty )}}=\exp \left[ -\Delta
p_{i+1}p_{i+1}-\frac{\Delta p_{i+3/2}-\Delta p_{i+1/2}}{4}\left(
p_{i+3/2}-p_{i+1/2}\right) \right]
\end{equation}

Using identities introduced in Sec. ***, 
\begin{equation}
\left\{ 
\begin{array}{l}
p_{i+3/2}=p_{i+1}+\frac{\Delta p_{i+3/2}}{2} \\ 
p_{i+1/2}=p_{i+1}-\frac{\Delta p_{i+1/2}}{2}
\end{array}
\right. 
\end{equation}
it is easy to demonstrate that 
\begin{equation}
p_{i+1}=\frac{1}{2}\left( p_{i+3/2}+p_{i+1/2}\right) -\frac{1}{2}\left( 
\frac{\Delta p_{i+3/2}}{2}-\frac{\Delta p_{i+1/2}}{2}\right) 
\end{equation}
and 
\begin{equation}
\Delta p_{i+1}=\frac{1}{2}\left( \Delta p_{i+3/2}+\Delta p_{i+1/2}\right)
=p_{i+3/2}-p_{i+1/2}
\end{equation}

Finally, since
\begin{eqnarray}
&&-\Delta p_{i+1}p_{i+1}-\frac{\Delta p_{i+3/2}-\Delta p_{i+1/2}}{4}\left(
p_{i+3/2}-p_{i+1/2}\right)   \nonumber \\
&=&-\left( p_{i+3/2}-p_{i+1/2}\right) \left[ \frac{1}{2}\left(
p_{i+3/2}+p_{i+1/2}\right) -\frac{1}{4}\left( \Delta p_{i+3/2}-\Delta %
p_{i+1/2}\right) \right]   \nonumber \\
&&-\frac{\Delta p_{i+3/2}-\Delta p_{i+1/2}}{4}\left(
p_{i+3/2}-p_{i+1/2}\right)   \nonumber \\
&=&-\frac{1}{2}\left( p_{i+3/2}+p_{i+1/2}\right) \left(
p_{i+3/2}-p_{i+1/2}\right)   \nonumber \\
&=&-\frac{1}{2}\left( p_{i+3/2}^{2}-p_{i+1/2}^{2}\right) 
\end{eqnarray}
one obtains the following relation 
\begin{equation}
\frac{f_{0,l+1/2,i+3/2,j+1/2}^{\left( 0\right) (\infty )}}{%
f_{0,l+1/2,i+1/2,j+1/2}^{\left( 0\right) (\infty )}}=\exp \left[ -\frac{%
p_{i+3/2}^{2}-p_{i+1/2}^{2}}{2}\right] 
\end{equation}
which is exactly corresponding to a numerical Maxwellian distribution
function. Hence, numerical errors do not propagate with the weighting here
considered.

Finally, for the non-uniform $p$ grid, the final results for interpolation
rules are
\begin{equation}
\delta _{p}^{\left( 0\right) }\rightarrow \left( 
\begin{array}{l}
\delta _{p,l+1/2,i+1,j+1/2}^{\left( 0\right) } \\ 
\delta _{p,l+1/2,i,j+1/2}^{\left( 0\right) }
\end{array}
\right) 
\end{equation}
with 
\begin{equation}
\delta _{p,l+1/2,i+1,j+1/2}^{\left( 0\right) }=\frac{1}{%
w_{p,l+1/2,i+1,j+1/2}^{\left( 0\right) (u)}}-\frac{1}{\exp \left[
w_{p,l+1/2,i+1,j+1/2}^{\left( 0\right) (u)}+w_{p,l+1/2,i+1,j+1/2}^{\left(
0\right) (nu)}\right] -1}
\end{equation}
and 
\begin{equation}
\delta _{p,l+1/2,i,j+1/2}^{\left( 0\right) }=\frac{1}{w_{p,l+1/2,i,j+1/2}^{%
\left( 0\right) (u)}}-\frac{1}{\exp \left[ w_{p,l+1/2,i,j+1/2}^{\left(
0\right) (u)}+w_{p,l+1/2,i,j+1/2}^{\left( 0\right) (nu)}\right] -1}
\end{equation}
and where 
\begin{equation}
w_{p}^{\left( 0\right) (u)}\rightarrow \left( 
\begin{array}{l}
w_{p,l+1/2,i+1,j+1/2}^{\left( 0\right) (u)} \\ 
w_{p,l+1/2,i,j+1/2}^{\left( 0\right) (u)}
\end{array}
\right) 
\end{equation}
with 
\begin{equation}
w_{p,l+1/2,i+1,j+1/2}^{\left( 0\right) (u)}=-\Delta p_{i+1}\frac{%
F_{p,l+1/2,i+1,j+1/2}^{\left( 0\right) }}{D_{pp,l+1/2,i+1,j+1/2}^{\left(
0\right) }}
\end{equation}
and 
\begin{equation}
w_{p,l+1/2,i,j+1/2}^{\left( 0\right) (u)}=-\Delta p_{i}\frac{%
F_{p,l+1/2,i,j+1/2}^{\left( 0\right) }}{D_{pp,l+1/2,i,j+1/2}^{\left(
0\right) }}
\end{equation}

Much in the same way, 
\begin{equation}
w_{p}^{\left( 0\right) (nu)}\rightarrow \left( 
\begin{array}{l}
w_{p,l+1/2,i+1,j+1/2}^{\left( 0\right) (nu)} \\ 
w_{p,i,j+1/2}^{\left( 0\right) (nu)}
\end{array}
\right) 
\end{equation}
where 
\begin{equation}
w_{p,l+1/2,i+1,j+1/2}^{\left( 0\right) (nu)}=-\frac{\Delta p_{i+3/2}-\Delta %
p_{i+1/2}}{4}\left[ \frac{F_{p,l+1/2,i+3/2,j+1/2}^{\left( 0\right) }}{%
D_{ppl+1/2,i+3/2,j+1/2}^{\left( 0\right) }}-\frac{F_{pl+1/2,i+1/2,j+1/2}^{%
\left( 0\right) }}{D_{ppl+1/2,i+1/2,j+1/2}^{\left( 0\right) }}\right] 
\end{equation}
and 
\begin{equation}
w_{p,i,j+1/2}^{\left( 0\right) (nu)}=-\frac{\Delta p_{i+1/2}-\Delta p_{i-1/2}%
}{4}\left[ \frac{F_{pl+1/2,i+1/2,j+1/2}^{\left( 0\right) }}{%
D_{ppl+1/2,i+1/2,j+1/2}^{\left( 0\right) }}-\frac{F_{pl+1/2,i-1/2,j+1/2}^{%
\left( 0\right) }}{D_{ppl+1/2,i-1/2,j+1/2}^{\left( 0\right) }}\right] 
\end{equation}

It is important to note that since only collisions are considered for
evaluating interpolating coefficients $\delta ^{\left( 0\right) },$ momentum
and pitch-angle dynamics are naturally decoupled. Consequently, the ratios $%
\frac{F_{p}^{\left( 0\right) }}{D_{pp}^{\left( 0\right) }}$ are only
functions of $p$, and it is the reason why they may take similar values at
different pitch-angle grid points $\xi _{0}.$

Calculations carried out for collisions only are usually extrapolated to
other accelaration mechanismes (Ohmic electric field, RF waves,...), using
the generalized weighting factor based on the simple rule, 
\begin{equation}
\left. \frac{F_{p}^{\left( 0\right) }}{D_{pp}^{\left( 0\right) }}\right|
_{effective}=\frac{\sum_{m}F_{p}^{\left( 0\right) (m)}}{\sum_{m}D_{pp}^{%
\left( 0\right) (m)}}
\end{equation}
where $\sum_{m}$ is the sum over all the physical processes $m$ that take
place in the plasma. Hence, the same syntax may be kept, 
\begin{equation}
w_{p,l+1/2,i+1,j+1/2}^{\left( 0\right) (nu)}=-\frac{\Delta p_{i+3/2}-\Delta 
p_{i+1/2}}{4}\left[ \left. \frac{F_{p}^{\left( 0\right) }}{D_{pp}^{\left(
0\right) }}\right| _{effective}^{l+1/2,i+3/2,j+1/2}-\left. \frac{%
F_{p}^{\left( 0\right) }}{D_{pp}^{\left( 0\right) }}\right|
_{effective}^{l+1/2,i+1/2,j+1/2}\right] 
\end{equation}
and 
\begin{equation}
w_{p,l+1/2,i,j+1/2}^{\left( 0\right) (nu)}=-\frac{\Delta p_{i+1/2}-\Delta %
p_{i-1/2}}{4}\left[ \left. \frac{F_{p}^{\left( 0\right) }}{D_{pp}^{\left(
0\right) }}\right| _{effective}^{l+1/2,i+1/2,j+1/2}-\left. \frac{%
F_{p}^{\left( 0\right) }}{D_{pp}^{\left( 0\right) }}\right|
_{effective}^{l+1/2,i-1/2,j+1/2}\right] 
\end{equation}

In that case, the interpolating weights exhibit a dependence with $\xi _{0}.$
It has however no importance for the accuracy of the calculations, since in
the domain where fluxes are strongly modified, the Maxwellian distribution
function is no more the solution of the Fokker-Planck equation. Usually, the
value of the Chang and Cooper function tends towards $\delta _{p}^{\left(
0\right) }\simeq \frac{1}{2}$ for a uniform mesh, which corresponds to a
standard linear interpolation.

Furthermore, $w_{p}^{\left( 0\right) (nu)}$ require evaluations of the
quantities at the grid points $\left( l+1/2,i-1/2,j+1/2\right) $ and $\left(
l+1/2,i+3/2,j+1/2\right) .$ The usual method to deal with this problem is
shifting indexes, according to the relation 
\begin{equation}
\left( l+1/2,i+1/2,j+1/2\right) \in \left\{ \left[ 1:n_{\psi }\right]
,\left[ 1:n_{p}\right] ,\left[ 1:n_{\xi }\right] \right\} 
\end{equation}
and 
\begin{equation}
\left\{ 
\begin{array}{l}
\left( l+1/2,i-1/2,j+1/2\right) \in \left\{ \left[ 1:n_{\psi }\right]
,\left[ X,1:n_{p}-1\right] ,\left[ 1:n_{\xi }\right] \right\}  \\ 
\left( l+1/2,i+3/2,j+1/2\right) \in \left\{ \left[ 1:n_{\psi }\right]
,\left[ 2:n_{p},X\right] ,\left[ 1:n_{\xi }\right] \right\} 
\end{array}
\right. 
\end{equation}

Hence, all quantities have the same size $\left( n_{\psi },n_{p},n_{\xi
}\right) ,$ which is crucial for the $3-D$ matrix representation.
Furthermore, boundary conditions are naturally satisfied with this
technique, since arbitrary values may be attributed to the variable $X$ at
these grid points.

\paragraph{Pitch-angle grid interpolation}

For the pitch-angle terms, calculation is very simple, and from previous
relations, one deduces directly that $w_{\xi ,l+1/2,i+1/2,j+1}^{\left(
0\right) }=w_{\xi ,l+1/2,i+1/2,j}^{\left( 0\right) }=0$ since $F_{\xi
}^{\left( 0\right) }=0$. Therefore, $\delta _{\xi ,l+1:2,i+1/2,j+1}^{\left(
0\right) }$ may be defined in an arbitrary manner. The most natural way is
to perform a linear interpolation, between $f_{0,l+1/2,i+1/2,j+3/2}^{\left(
0\right) (\infty )}$ and $f_{0,l+1/2,i+1/2,j+3/2}^{\left( 0\right) (\infty
)},$ and also between $f_{0,l+1:2,i+1/2,j+1/2}^{\left( 0\right) (\infty )}$
and $f_{0,l+1/2,i+1/2,j-1/2}^{\left( 0\right) (\infty )}$ according to the
relations 
\begin{eqnarray}
f_{0,l+1/2,i+1/2,j+1}^{\left( 0\right) (\infty )} &=&\left( \frac{\Delta \xi
_{0,j+1/2}}{\Delta \xi _{0,j+1/2}+\Delta \xi _{0,j+3/2}}\right)
f_{0,l+1/2,i+1/2,j+3/2}^{\left( 0\right) (\infty )}  \nonumber \\
&&+\left( \frac{\Delta \xi _{0,j+3/2}}{\Delta \xi _{0,j+1/2}+\Delta \xi
_{0,j+3/2}}\right) f_{0,l+1/2,i+1/2,j+1/2}^{\left( 0\right) (\infty )}
\end{eqnarray}
\begin{eqnarray}
f_{0,l+1/2,i+1/2,j}^{\left( 0\right) (\infty )} &=&\left( \frac{\Delta \xi
_{0,j-1/2}}{\Delta \xi _{0,j-1/2}+\Delta \xi _{0,j+1/2}}\right)
f_{0,l+1/2,i+1/2,j+1/2}^{\left( 0\right) (\infty )}  \nonumber \\
&&+\left( \frac{\Delta \xi _{0,j+1/2}}{\Delta \xi _{0,j-1/2}+\Delta \xi
_{0,j+1/2}}\right) f_{0,l+1/2,i+1/2,j-1/2}^{\left( 0\right) (\infty )}
\end{eqnarray}

Taking into account of the non-uniform grid steps 
\begin{equation}
\delta _{\xi ,l+1/2,i+1/2,j+1}^{\left( 0\right) }=\frac{\Delta \xi _{0,j+3/2}%
}{\Delta \xi _{0,j+1/2}+\Delta \xi _{0,j+3/2}}
\end{equation}
\begin{equation}
\delta _{\xi ,l+1/2,i+1/2,j}^{\left( 0\right) }=\frac{\Delta \xi _{0,j+1/2}}{%
\Delta \xi _{0,j-1/2}+\Delta \xi _{0,j+1/2}}
\end{equation}
and for a uniform grid, the relation $\delta _{\xi ,l+1/2,i+1/2,j+1}^{\left(
0\right) }=\delta _{\xi ,l+1/2,i+1/2,j}^{\left( 0\right) }=1/2$ are exactly
recovered.

Finally, for the non-uniform pitch-angle grid $\xi ,$ the interpolation
rules are
\begin{equation}
\delta _{\xi }^{\left( 0\right) }\rightarrow \left\{ 
\begin{array}{l}
\delta _{\xi ,l+1/2,i+1/2,j+1}^{\left( 0\right) }=\frac{\Delta \xi _{0,j+3/2}%
}{\Delta \xi _{0,j+1/2}+\Delta \xi _{0,j+3/2}} \\ 
\delta _{\xi ,l+1/2,i+1/2,j}^{\left( 0\right) }=\frac{\Delta \xi _{0,j+1/2}}{%
\Delta \xi _{0,j-1/2}+\Delta \xi _{0,j+1/2}}
\end{array}
\right. 
\end{equation}

Furthermore, $\delta _{\xi }^{\left( 0\right) }$ require evaluations of the
quantities at the grid points $\left( l+1/2,i+1/2,j-1/2\right) $ and $\left(
l+1/2,j+1/2,j+3/2\right) $. As for the momentum grid interpolation, indexes
are shifted, 
\begin{equation}
\left\{ 
\begin{array}{l}
\left( l+1/2,i+1/2,j-1/2\right) \in \left\{ \left[ 1:n_{\psi }\right]
,\left[ 1:n_{p}\right] ,\left[ X,1:n_{\xi }-1\right] \right\}  \\ 
\left( l+1/2,i+1/2,j+3/2\right) \in \left\{ \left[ 1:n_{\psi }\right]
,\left[ 1:n_{p}\right] ,\left[ 2:n_{\xi },X\right] \right\} 
\end{array}
\right. 
\end{equation}

\paragraph{Spatial grid interpolation}

The determination of $\delta _{\psi }^{\left( 0\right) }$ is based on the
same method used for $\delta _{\xi }^{\left( 0\right) }$, since no specific
condition is required for interpolating  $f_{0}^{\left( 0\right) }$ on the
flux grid. A linear interpolation is considered between $%
f_{0,l+3/2,i+1/2,j+1/2}^{\left( 0\right) (\infty )}$ and $%
f_{0,l+1/2,i+1/2,j+1/2}^{\left( 0\right) (\infty )}$ and also between $%
f_{0,l+1/2,i+1/2,j+1/2}^{\left( 0\right) (\infty )}$ and $%
f_{0,l-1/2,i+1/2,j+1/2}^{\left( 0\right) (\infty )}$, 
\begin{eqnarray}
f_{0,l+1,i+1/2,j+1/2}^{\left( 0\right) (\infty )} &=&\left( \frac{\Delta %
\psi _{l+1/2}}{\Delta \psi _{l+1/2}+\Delta \psi _{l+3/2}}\right)
f_{0,l+3/2,i+1/2,j+1/2}^{\left( 0\right) (\infty )}  \nonumber \\
&&+\left( \frac{\Delta \psi _{l+3/2}}{\Delta \psi _{l+1/2}+\Delta \psi
_{l+3/2}}\right) f_{0,l+1/2,i+1/2,j+1/2}^{\left( 0\right) (\infty )}
\end{eqnarray}

\begin{eqnarray}
f_{0,l,i+1/2,j+1/2}^{\left( 0\right) (\infty )} &=&\left( \frac{\Delta \psi
_{l-1/2}}{\Delta \psi _{l-1/2}+\Delta \psi _{l+1/2}}\right)
f_{0,l+1/2,i+1/2,j+1/2}^{\left( 0\right) (\infty )}  \nonumber \\
&&+\left( \frac{\Delta \psi _{l+1/2}}{\Delta \psi _{l-1/2}+\Delta \psi
_{l+1/2}}\right) f_{0,l-1/2,i+1/2,j+1/2}^{\left( 0\right) (\infty )}
\end{eqnarray}
taking into account of the non-uniform nature of the grid 
\begin{equation}
\delta _{\psi ,l+1}^{\left( 0\right) }=\frac{\Delta \psi _{l+3/2}}{\Delta %
\psi _{l+1/2}+\Delta \psi _{l+3/2}}
\end{equation}
\begin{equation}
\delta _{\psi ,l}^{\left( 0\right) }=\frac{\Delta \psi _{l+1/2}}{\Delta \psi
_{l-1/2}+\Delta \psi _{l+1/2}}
\end{equation}
For a uniform grid the relation $\delta _{\psi ,l+1}^{\left( 0\right)
}=\delta _{\psi ,l}^{\left( 0\right) }=1/2$ is well recovered.

Furthermore, $\delta _{\psi }^{\left( 0\right) }$ require evaluations of the
quantities at the grid points $\left( l-1/2,i+1/2,j+1/2\right) $ and $\left(
l+3/2,j+1/2,j+1/2\right) $. As for the momentum grid interpolation, indexes
are shifted according to the rule 
\begin{equation}
\left\{ 
\begin{array}{l}
\left( l-1/2,i+1/2,j+1/2\right) \in \left\{ \left[ X,1:n_{\psi }-1\right]
,\left[ 1:n_{p}\right] ,\left[ 1:n_{\xi }\right] \right\}  \\ 
\left( l+3/2,i+1/2,j+1/2\right) \in \left\{ \left[ 2:n_{\psi },X\right]
,\left[ 1:n_{p}\right] ,\left[ 1:n_{\xi }\right] \right\} 
\end{array}
\right. 
\end{equation}
